% problem-type: poisson-neumann-tuong-thich
% order: 90
% Nguon: De thi MAT3365 (PTDHR), HK2 2023-2024, K67 Toan tin, Cau 1. (tu giai, KHONG co dap an chinh thuc)
% Bo cuc: De vi du -> Cong thuc -> De + Bai giai

\section{Poisson--Neumann trong quạt --- điều kiện tương thích \& năng lượng}

% ---------------------------------------------------------------
\subsection*{Đề bài ví dụ}
% ---------------------------------------------------------------
Xét bài toán biên Neumann cho phương trình Poisson
\[
  \Delta u=1\ \text{trong } Q=\{(x,y): 1<y<x,\ (x-1)^2+(y-1)^2<1\},
\]
với $\partial_\nu u=0$ khi $(x-y)(1-y)=0$ (hai cạnh thẳng) và $\partial_\nu u=x+C$ trên cung
$(x-1)^2+(y-1)^2=1$. \ (a) Tìm $C$ để có nghiệm. (b) Bằng tích phân năng lượng, chỉ ra có vô số
nghiệm sai khác hằng số. (c) Giải bài toán.

\emph{\small(Nguồn: Đề thi MAT3365 --- PTĐHR, HK2 2023--2024, K67 Toán tin, Câu 1. Lời giải dưới đây tự soạn.)}

% ---------------------------------------------------------------
\subsection*{Nhận ra dạng này khi nào?}
% ---------------------------------------------------------------
\begin{itemize}
  \item Bài toán \textbf{Neumann} (cho $\partial_\nu u$ trên TOÀN biên) cho Laplace/Poisson.
  \item Phải kiểm tra \textbf{điều kiện tương thích} mới có nghiệm; nghiệm KHÔNG duy nhất (sai khác hằng).
\end{itemize}
\begin{meobox}
Tương thích: tổng thông lượng biên $\int_{\partial}\partial_\nu u=$ tích phân nguồn $\iint_Q\Delta u$ ($=$ diện tích nếu nguồn $=1$).
\end{meobox}

% ---------------------------------------------------------------
\subsection*{Công thức \& phương pháp}
% ---------------------------------------------------------------
\subsubsection*{Tích phân năng lượng (sai khác hằng)}
% technique: tich-phan-nang-luong
\textbf{Tích phân năng lượng} (chứng minh DUY NHẤT nghiệm / nghiệm sai khác hằng).
Xét hiệu $u=u_1-u_2$ của hai nghiệm $\Rightarrow$ $u$ thỏa bài toán \emph{thuần nhất} (dữ liệu $=0$).
\begin{itemize}
  \item \emph{Elliptic} ($\Delta u=0$, biên $u=0$ hoặc $\partial_\nu u=0$): Green $\Rightarrow$
  $\iint_\Omega|\nabla u|^2=\int_{\partial\Omega}u\,\partial_\nu u=0\Rightarrow \nabla u=0\Rightarrow u$ hằng (Dirichlet: $u=0$).
  \item \emph{Parabolic} ($u_t=k\Delta u$): đặt $I(t)=\int u^2\,dx\ge0$; $I'(t)=-2k\int|\nabla u|^2\le0$ và $I(0)=0\Rightarrow I\equiv0\Rightarrow u=0$.
\end{itemize}

\subsubsection*{Tách biến cực để giải}
% method: tach-bien-cuc
\begin{dieukienbox}
Dùng khi: $\Delta u=0$ (Laplace) hoặc $\Delta u=$ hằng (Poisson) trên \textbf{quạt / vành khăn / đĩa},
biên gồm cung tròn ($r=$ const) và cạnh thẳng ($\theta=$ const).
\end{dieukienbox}

\begin{nhobox}
\textbf{Laplacian cực:} $\Delta u=u_{rr}+\frac1r u_r+\frac{1}{r^2}u_{\theta\theta}$. Tách $u=R(r)\Theta(\theta)$:
\begin{itemize}
  \item Góc: $\Theta''+\mu^2\Theta=0\Rightarrow \Theta=\cos(\mu\theta)$ hoặc $\sin(\mu\theta)$ ($\mu$ từ biên ở 2 cạnh $\theta$).
  \item Bán kính (Euler): $r^2R''+rR'-\mu^2 R=0\Rightarrow R=r^{\mu},\,r^{-\mu}$ (hoặc $a+b\ln r$ khi $\mu=0$).
\end{itemize}
Nghiệm tổng quát:
$\displaystyle u=a_0+b_0\ln r+\sum_{n\ge1}\big(a_n r^{\mu_n}+b_n r^{-\mu_n}\big)\Theta_n(\theta).$
\end{nhobox}

\begin{meobox}
\begin{itemize}
  \item Cạnh Neumann $u_\theta=0$ tại $\theta=0$ $\Rightarrow$ chọn $\cos(\mu\theta)$; Dirichlet $u=0$ $\Rightarrow$ chọn $\sin(\mu\theta)$. Cạnh còn lại định $\mu_n$.
  \item \textbf{Đĩa} (chứa gốc): bỏ $r^{-\mu},\ln r$ (chặn tại $0$). \textbf{Vành khăn}: giữ cả hai.
  \item Poisson $\Delta u=1$: cộng nghiệm riêng $r^2/4$. Hệ số $a_n,b_n$ lấy từ khai triển Fourier dữ liệu trên cung.
\end{itemize}
\end{meobox}
\emph{(Gợi ý đề: đổi sang cực tâm tại đỉnh quạt $x=1+r\cos\theta,\ y=1+r\sin\theta$.)}

% ---------------------------------------------------------------
\subsection*{Đề bài \& bài giải}
% ---------------------------------------------------------------
\paragraph{Hình học miền.} Đổi $x=1+r\cos\theta,\ y=1+r\sin\theta$ (tâm tại đỉnh $(1,1)$): điều kiện
$y>1\Rightarrow\sin\theta>0$; $y<x\Rightarrow\sin\theta<\cos\theta\Rightarrow\theta<\tfrac{\pi}{4}$. Vậy $Q$ là
\textbf{quạt} $0<\theta<\tfrac{\pi}{4},\ 0<r<1$; hai cạnh thẳng là $\theta=0$ và $\theta=\tfrac{\pi}{4}$, cung là $r=1$.

\paragraph{(a) Điều kiện tương thích.} Green: $\displaystyle\int_{\partial Q}\partial_\nu u\,ds=\iint_Q\Delta u\,dA=|Q|$.
Diện tích quạt $|Q|=\tfrac12 r^2\Delta\theta=\tfrac12\cdot1\cdot\tfrac{\pi}{4}=\tfrac{\pi}{8}$. Hai cạnh thẳng cho $0$;
trên cung $r=1$ có $x=1+\cos\theta$, $ds=d\theta$:
\[
  \int_0^{\pi/4}(1+\cos\theta+C)\,d\theta=(1+C)\tfrac{\pi}{4}+\tfrac{\sqrt2}{2}.
\]
Cân bằng với $\tfrac{\pi}{8}$: $(1+C)\tfrac{\pi}{4}=\tfrac{\pi}{8}-\tfrac{\sqrt2}{2}\Rightarrow
\boxed{C=-\tfrac12-\tfrac{2\sqrt2}{\pi}}$.

\paragraph{(b) Sai khác hằng (năng lượng).} Hai nghiệm $u_1,u_2$; hiệu $w=u_1-u_2$ thỏa $\Delta w=0$,
$\partial_\nu w=0$. Green: $\iint_Q|\nabla w|^2=\int_{\partial Q}w\,\partial_\nu w=0\Rightarrow\nabla w=0\Rightarrow w$ hằng.
Vậy nghiệm tồn tại thì có vô số, sai khác một hằng số cộng.

\paragraph{(c) Bước 1 --- nghiệm tổng quát.} $\Delta u=1$ có nghiệm riêng $\tfrac{r^2}{4}$ ($\Delta\tfrac{r^2}{4}=1$,
lại $u_\theta=0$ nên thỏa Neumann hai cạnh). Phần điều hòa với $u_\theta(r,0)=0$ chọn $\cos(\mu\theta)$;
$u_\theta(r,\tfrac{\pi}{4})=0\Rightarrow\mu=4n$. Bỏ $r^{-4n}$ (chặn tại $0$):
\[
  u(r,\theta)=\frac{r^2}{4}+A_0+\sum_{n\ge1}A_n\,r^{4n}\cos(4n\theta).
\]

\paragraph{(c) Bước 2 --- khớp cung $r=1$.} $\partial_r u|_{r=1}=\tfrac12+\sum_{n\ge1}4nA_n\cos(4n\theta)$
phải bằng $x+C=1+\cos\theta+C$. Chuyển vế ($\tfrac12+C=-\tfrac{2\sqrt2}{\pi}$):
\[
  \sum_{n\ge1}4nA_n\cos(4n\theta)=\cos\theta-\tfrac{2\sqrt2}{\pi}.
\]
Hệ $\{\cos(4n\theta)\}$ trực giao trên $[0,\tfrac{\pi}{4}]$ ($\int\cos^2=\tfrac{\pi}{8}$); số hạng hằng tự khử
(đúng điều kiện (a)). Với $n\ge1$:
\[
  4nA_n=\tfrac{8}{\pi}\!\int_0^{\pi/4}\!\cos\theta\cos(4n\theta)\,d\theta
  =\tfrac{8}{\pi}\cdot\frac{(-1)^{n+1}\sqrt2}{2(16n^2-1)}
  \Rightarrow A_n=\frac{(-1)^{n+1}\sqrt2}{n\pi(16n^2-1)}.
\]
\begin{nhobox}
\[
  u(r,\theta)=\frac{r^2}{4}+A_0+\frac{\sqrt2}{\pi}\sum_{n\ge1}\frac{(-1)^{n+1}}{n(16n^2-1)}\,r^{4n}\cos(4n\theta),
\]
với $A_0$ tùy ý (sai khác hằng) và $C=-\tfrac12-\tfrac{2\sqrt2}{\pi}$; $x=1+r\cos\theta,\ y=1+r\sin\theta$.
\end{nhobox}
